\subsection{Operational room allocation and deterministic decoding}
\label{sec:room-formulation}

At decision time $t$, let $W$ be the waiting cases, $\mathcal R$ the rooms,
$N=|W|$, $[a,b)$ the opening window, $\tau$ the turnover time and $A_r$
the estimated availability of room $r$. Case $i$ has integer predicted
occupancy $p_i>0$. Running and completed assignments form a fixed set $F$.
The search chooses binary allocations $y_{ir}$ satisfying
\begin{equation}
 \sum_{r\in\mathcal R}y_{ir}=1\quad(i\in W),\qquad y_{ir}\in\{0,1\}.
\end{equation}
All rooms are eligible. Allocation does not guarantee that a case can start:
the decoder may return it explicitly unassigned.

Let $h_r(v)$ return the first minute at or after $v$ outside any currently
known start-blocking closure of room $r$. For each room, order its allocated
cases by $(p_i,\text{local ID}_i)$. Initialise $v_r=\max(t,a,A_r)$, then
compute $s_i=h_r(v_r)$. If $s_i<b$, accept the assignment and update
$v_r=s_i+p_i+\tau$; otherwise mark the case unassigned without advancing
$v_r$. Restore $F$ unchanged. This recurrence fixes the sequencing rule and
ensures predicted non-overlap with turnover; it does not optimise permutations.
Starts must satisfy $s_i\geq\max(t,a,A_r)$ and avoid known closures. Finishes
may exceed $b$. Future closures are absent until revealed by execution.

Write $D(y)$ for the decoded schedule and $W_A(y)$ for its assigned waiting
cases. Let $q_i(y)$ denote the complete room/start/duration assignment,
or the unassigned symbol, and let $q_i^0$ denote the accepted decision.
The request-level components are
\begin{align}
 U(y)&=N-|W_A(y)|,\\
 O(y)&=\sum_{i\in W_A(y)}[s_i+p_i+\tau-b]_+,\\
 R(y)&=\sum_{i\in W}\mathbf{1}\{q_i(y)\ne q_i^0\},
 \qquad B=\sum_{i\in W}(p_i+\tau).
\end{align}
At initial planning, set $R=0$. The implemented minimisation objective is
\begin{equation}
 C_{\mathrm{room}}(y)=U(y)+\frac{O(y)+R(y)/(N+1)}{B+1},
 \qquad f(y)=-C_{\mathrm{room}}(y).
 \label{eq:room-cost}
\end{equation}
Since starts precede closing, $O\leq B$; also $R\leq N$. The fractional
term is less than one, so an additional unstarted case dominates every
possible overtime/change improvement. Overtime uses integer minutes, and
$R/(N+1)<1$, so one minute dominates all changes. This encodes the intended
lexicographic priorities without clinically calibrated trade-off weights.
Fixed activity's overtime is constant within a request and excluded from
this score; execution metrics include it.

The admissible schedules must cover each case exactly once as assigned or
unassigned, retain $F$, and pass the common temporal validator. Enumeration
on a small instance certifies the best cost only among allocations decoded
by this rule. It does not certify optimal execution with hidden durations,
an optimal independent sequence, or an optimum across objectives. Different
workloads or duration modes change $B$, so raw normalised costs are not a
cross-instance performance scale.
